\jabstract{
本研究は、医学生物学分野における高被引用論文の長期的な引用ダイナミクスおよび構造的特徴を明らかにすることを目的とする。従来の計量書誌学では、3年・5年インパクトファクターのような短期的引用指標が主に用いられてきたが、論文の真の影響を把握するには長期的観察が不可欠である。本研究では、単なる被引用数ではなく、引用寿命、雑誌構造、研究テーマ構造といった多層的観点から分析を行った。
データにはPubMed annual baseline (2023) を用い、1946年から2020年までに出版された約3000万件の論文に記載されるレファレンス情報を対象とした。これに対して5年間ごとの「累積世代モデル」を構築し、各累積世代の被引用数全体の1\%を構成する上位の論文を「高被引用論文」と定義した。また、高被引用論文とそれ以外の論文との比較のために4分位グループを設定した。
論文単位の分析の結果、高被引用論文の中でも、最多被引用論文については数十年にわたり引用が持続し現在までオブソレッセンス（被引用減衰）に至らない長寿命型の論文が存在することが確認された。一方、高被引用論文の中でも被引用減衰が確認された論文では、出版から被引用ピークまで約4年、ピークから減衰まで約7年という寿命構造が観察され、生物分野における既存モデルと整合的であった。被引用減衰が確認された論文の階層クラスタリングの結果、明確なタイプ分化は確認されなかった。
雑誌レベルの分析の結果、高被引用論文を多く掲載する雑誌が必ずしも高インパクトファクター誌であるとは限らないことが示された。ただし、Kullback-Leibler divergence による分析から、高被引用論文は他グループと比較して特定雑誌への集中傾向（嗜好）が強いことが明らかとなった。
研究テーマ構造の分析では、HeSHディスクリプターを使ったエントロピーの分析から、高被引用論文は他グループと比較して研究テーマが集中的であり、それを引用する論文の研究テーマもまた集中的であることが明らかとなった。一方で研究テーマとそれを引用する論文の研究テーマの間には乖離が見られ、他のグループに比べ異なる分野から引用されやすい傾向にあることが明らかになった。
}

\eabstract{
This study aims to elucidate the long-term citation dynamics and structural characteristics of highly cited papers in the field of biomedical sciences. In conventional bibliometrics, short-term citation indicators such as the 3-year and 5-year impact factors have been predominantly used; however, long-term observation is essential for accurately capturing the true impact of scholarly articles. In this study, analysis was conducted from multiple perspectives, including citation lifespan, journal structure, and research topic structure, rather than relying solely on citation counts.
The dataset used was the PubMed Annual Baseline (2023), covering reference information from approximately 30 million papers published between 1946 and 2020. Based on this dataset, a “cumulative generation model” was constructed at five-year intervals, and the top 1\% of papers in terms of total citations within each cumulative generation were defined as “highly cited papers.” In addition, quartile groups were established to enable comparisons between highly cited papers and other papers.
At the individual paper level, the analysis revealed that even among highly cited papers, the most cited ones exhibit sustained citation over several decades, forming a class of long-lived papers that have not yet reached obsolescence (citation decay). Conversely, among highly cited papers where citation decay was observed, a lifecycle structure was identified: approximately four years from publication to peak citation, and about seven years from the peak to decline, which is consistent with existing models in the biological sciences. Hierarchical clustering of papers exhibiting citation decay did not reveal clear typological differentiation.
At the journal level, it was shown that journals publishing many highly cited papers are not necessarily those with high impact factors. However, analysis using Kullback–Leibler divergence demonstrated that highly cited papers tend to exhibit a stronger concentration (preference) toward specific journals compared to other groups.
At the research topic level, entropy analysis based on MeSH descriptors revealed that highly cited papers are more concentrated in their research topics than other groups, and the topics of papers citing them are also similarly concentrated. At the same time, a divergence was observed between the topics of the cited papers and those of the citing papers, indicating that highly cited papers tend to be cited across different research fields.
}

